%% OVERFORE-Subdocumento8-Anexos.tex — generado desde OVERFORE-Subdocumento8-Anexos.docx
%% Formato: clase overfore, modo anexo. El formato lo pone la clase; aquí sólo va el contenido.
\documentclass[anexo]{../../comun/overfore}
\numsubdoc{8}
\informe{Informe Preparatorio N° 2}
\fecha{12 de octubre de 2026}
\externos{S5,S8,T14}
\begin{document}

\anexo{Escalas, umbrales y reglas de escalamiento}
Este anexo detalla las escalas que las secciones 8.1.2 a 8.1.4 del Subdocumento 8 resumen. Muestra de dónde sale cada umbral de impacto, cuáles son los siete umbrales que ordenan el escalamiento y cómo se declara el apetito de riesgo de Overfore para este contrato. Todo impacto se expresa en semanas, horas hombre (HH) o porcentaje del hito, nunca en montos, conforme al Art. 50.2 de las Bases administrativas (p. 29).

\seccion{Derivación de los umbrales de esfuerzo}
Los umbrales de esfuerzo se fijan como una fracción del esfuerzo de las ramas 1.1 a 1.7, para que un mismo nivel signifique lo mismo en cualquier riesgo. La base son las 61.194 HH de las ramas 1.1 a 1.7 del Formulario T-14 (Tabla \refext{T14}{tab:esfuerzo-estimado-rama-origen-horas} del Formulario T-14), que ya son públicas en la oferta técnica y no permiten inferir su monto. La Tabla \ref{tab:derivacion-umbrales-esfuerzo} muestra el cálculo.

\begin{tabla}{I{0.171}D{0.189}I{0.304}D{0.336}}{Derivación de los umbrales de esfuerzo.}{tab:derivacion-umbrales-esfuerzo}
  {\cab{Nivel de impacto} & \cab{Fracción de referencia} & \cab{Cálculo} & \cab{Umbral usado}}
  {Fuente: elaboración propia a partir de Formulario T-14 (asociado al cap. 7), tabla \refext{T14}{tab:esfuerzo-estimado-rama-origen-horas} del Formulario T-14 (esfuerzo estimado por rama).}
Muy alto & 5 \% & 0,05 × 61.194 = 3.060 HH & Más de 3.000 HH \\
Alto & 2 \% & 0,02 × 61.194 = 1.224 HH & 1.200 a 3.000 HH, inclusive \\
Medio & 1 \% & 0,01 × 61.194 = 612 HH & 600 o más y menos de 1.200 HH \\
Bajo & 0,5 \% & 0,005 × 61.194 = 306 HH & 300 o más y menos de 600 HH \\
Muy bajo & Menos de 0,5 \% & No aplica & Menos de 300 HH \\
\end{tabla}
De la tabla se desprende que los umbrales se redondean a la centena inferior para que sean fáciles de recordar en un comité, sin cambiar el nivel de ningún riesgo del registro.

\seccion{Derivación de los umbrales de plazo}
Los umbrales de plazo se miden como atraso bruto sobre una cadena sin holgura del Formulario T-14 (Tabla \refext{T14}{tab:red-precedencias-calculo-ruta-critica-meses} del Formulario T-14), antes de usar la reserva de plazo. Se anclan en la multa por atraso del Art. 80, de 0,5 \% del valor del hito por día corrido con tope de 10 \%, y en la causal de término del Art. 81, como muestra la Tabla \ref{tab:derivacion-umbrales-plazo}.

\begin{tabla}{I{0.174}D{0.211}D{0.236}I{0.379}}{Derivación de los umbrales de plazo.}{tab:derivacion-umbrales-plazo}
  {\cab{Nivel de impacto} & \cab{Atraso} & \cab{Multa equivalente sobre el hito} & \cab{Fundamento}}
  {Fuente: elaboración propia a partir de las Bases administrativas (Compañía Minera Altos de Aranda [CMAA], 2026a, cap. 21, Art. 80, pp. 41–42); Bases administrativas, cap. 22, art. 81, p. 42.}
Muy alto & Más de 30 días & Tope alcanzado & Más de 30 días corridos abre una causal de término \\
Alto & 20 a 30 días, inclusive & 10 \% & A los 20 días corridos la multa llega a su tope \\
Medio & 7 o más y menos de 20 días & 3,5 \% a menos de 10 \% & 7 o más y menos de 20 días corridos \\
Bajo & 2 o más y menos de 7 días & 1 \% a menos de 3,5 \% & 2 o más y menos de 7 días corridos \\
Muy bajo & Menos de 2 días & Menos de 1 \% & Absorbible en una ventana de paso \\
\end{tabla}
El nivel alto se ancla en 20 días corridos, cuando se alcanza el tope acumulado de la multa; el incumplimiento existe desde el primer día de atraso contractual. Para estimar esfuerzo se emplean dos equivalencias del T-14: una semana de construcción de la cadena crítica E1 equivale a 342 HH, la mitad de un sprint de 684 HH; una semana media de marcha blanca E1 equivale a 565 HH, redondeando 40 HH de 1.6.3.1 más 6.824 / 13 = 524,923 HH de 1.6.3.3. R-12 conserva 1.695 HH inherentes y 1.130 HH residuales. R-23 conserva 565 HH por una semana; su impacto alto deriva de la extensión del cierre, no de aumentar estas horas. La cobertura efectiva se dimensiona por áreas y turnos.

\seccion{Umbrales de escalamiento}
El escalamiento no depende del criterio de quien preside el comité, sino de siete umbrales numerados. La Tabla \ref{tab:umbrales-escalamiento-plan-riesgos} los presenta con la acción que gatillan.

\begin{tabla}{D{0.062}I{0.435}I{0.503}}{Umbrales de escalamiento del plan de riesgos.}{tab:umbrales-escalamiento-plan-riesgos}
  {\cab{N°} & \cab{Umbral} & \cab{Acción}}
  {Fuente: elaboración propia a partir de Formulario T-14 (asociado al cap. 7), sección 7.}
1 & Riesgo nuevo o recalificado con exposición de 0,40 o más & Comité Ejecutivo extraordinario en 5 días hábiles con propuesta de cambio de diseño, de plan o de alcance \\
2 & Riesgo con exposición de 0,20 o más sin respuesta aprobada & Plan de respuesta al Comité Ejecutivo siguiente \\
3 & Consumo de 50 \% o más de una reserva de contingencia & Plan de recuperación en el Comité de Proyecto \\
4 & Consumo de 80 \% o más de una reserva, o de 3,4 semanas de la reserva de plazo del mes 11 & Escala al Comité Ejecutivo, único que autoriza la reserva de gestión \\
5 & Probabilidad de superar la reserva de plazo del mes 11, recalculada, mayor que 20 \% & Comité Ejecutivo y aviso al cliente si el hito se desvía más de 10 \% \\
6 & Índice de desempeño del cronograma entre 0,90 y 0,95, o bajo 0,90 & Revisión en el Comité de Proyecto, o plan de recuperación \\
7 & Disparador activado de un riesgo exigido por las Bases técnicas del caso Minería & Se informa en el Comité Ejecutivo del mes, aunque su exposición sea media \\
\end{tabla}
De la tabla se desprende que los umbrales 3 y 4 equivalen a 2,2 y 3,4 semanas de la reserva de plazo del mes 11, la mitad y el 80 \% de sus 4,3 semanas. El Anexo D calcula la probabilidad de alcanzarlos.

\seccion{Apetito, tolerancia y umbral}
El apetito de Overfore es bajo en plazo y en seguridad: la oferta se presenta sólo si ningún riesgo residual queda en zona alta o crítica. Es moderado en innovación, donde acepta riesgos de adopción en zona media con verificación periódica. La tolerancia es de cero días de atraso imputable en los hitos H5, H7 y H12, una reserva de contingencia de plazo de la Etapa 1 no mayor que la mitad de la reserva del mes 11 y ninguna vulnerabilidad crítica abierta más de 7 días corridos. El umbral está formado por los siete escalamientos de la Tabla \ref{tab:umbrales-escalamiento-plan-riesgos}, y la reserva de gestión sólo la libera el Comité Ejecutivo.

\anexo{Cálculo de la probabilidad, el impacto y la exposición}
Este anexo contiene el cálculo que las secciones 8.2.3 y 8.3.1 del Subdocumento 8 resumen en las Figuras \refext{S8}{fig:ubicacion-24-riesgos-matriz-exposicion-antes} y \refext{S8}{fig:ubicacion-24-riesgos-matriz-exposicion-despues}. Para cada riesgo, la exposición es E = P × I, con P en su estimación puntual e I en el valor de su nivel (0,1; 0,3; 0,5; 0,7 o 0,9). El valor esperado de plazo es P por las semanas brutas sobre su cadena o tramo, y el de esfuerzo es P por las horas hombre adicionales. El enunciado y el fundamento de cada estimación están en la ficha del riesgo en el Formulario T-16.

\seccion{Exposición inherente}
La Tabla \ref{tab:exposicion-valor-esperado-inherentes-24-riesgos} presenta los 24 riesgos antes de aplicar las respuestas, con la zona que resulta de su exposición.

\begin{tabla}{I{0.081}D{0.137}D{0.062}D{0.118}I{0.105}D{0.130}D{0.099}D{0.130}D{0.137}}{Exposición y valor esperado inherentes de los 24 riesgos.}{tab:exposicion-valor-esperado-inherentes-24-riesgos}
  {\cab{ID} & \cab{P} & \cab{I} & \cab{E = P × I} & \cab{Zona} & \cab{Semanas} & \cab{HH} & \cab{VE plazo (sem)} & \cab{VE esfuerzo (HH)}}
  {Fuente: elaboración propia a partir de Formulario T-16 (asociado al cap. 8), sección 2.}
R-01 & 0,45 & 0,5 & 0,225 & Alta & 1,5 & 500 & 0,68 & 225 \\
R-02 & 0,15 & 0,5 & 0,075 & Baja & 0,0 & 800 & 0,00 & 120 \\
R-03 & 0,15 & 0,5 & 0,075 & Baja & 1,0 & 400 & 0,15 & 60 \\
R-04 & 0,20 & 0,9 & 0,180 & Media & 0,0 & 300 & 0,00 & 60 \\
R-05 & 0,60 & 0,7 & 0,420 & Crítica & 2,5 & 400 & 1,50 & 240 \\
R-06 & 0,15 & 0,7 & 0,105 & Media & 3,0 & 480 & 0,45 & 72 \\
R-07 & 0,45 & 0,7 & 0,315 & Alta & 1,5 & 400 & 0,68 & 180 \\
R-08 & 0,60 & 0,7 & 0,420 & Crítica & 2,0 & 1.130 & 1,20 & 678 \\
R-09 & 0,25 & 0,7 & 0,175 & Media & 2,0 & 200 & 0,50 & 50 \\
R-10 & 0,35 & 0,5 & 0,175 & Media & 1,5 & 310 & 0,53 & 109 \\
R-11 & 0,55 & 0,5 & 0,275 & Alta & 0,0 & 384 & 0,00 & 211 \\
R-12 & 0,60 & 0,7 & 0,420 & Crítica & 3,0 & 1.695 & 1,80 & 1.017 \\
R-13 & 0,30 & 0,5 & 0,150 & Media & 0,0 & 160 & 0,00 & 48 \\
R-14 & 0,35 & 0,5 & 0,175 & Media & 0,0 & 160 & 0,00 & 56 \\
R-15 & 0,4557375 & 0,5 & 0,227869 & Alta & 2,0 & 684 & 0,911475 & 311,724 \\
R-16 & 0,50 & 0,7 & 0,350 & Alta & 0,0 & 400 & 0,00 & 200 \\
R-17 & 0,25 & 0,7 & 0,175 & Media & 0,0 & 384 & 0,00 & 96 \\
R-18 & 0,25 & 0,7 & 0,175 & Media & 2,0 & 240 & 0,50 & 60 \\
R-19 & 0,30 & 0,7 & 0,210 & Alta & 3,0 & 1.026 & 0,90 & 308 \\
R-20 & 0,30 & 0,7 & 0,210 & Alta & 0,0 & 400 & 0,00 & 120 \\
R-21 & 0,40 & 0,5 & 0,200 & Alta & 0,0 & 200 & 0,00 & 80 \\
R-22 & 0,10 & 0,7 & 0,070 & Baja & 0,0 & 800 & 0,00 & 80 \\
R-23 & 0,35 & 0,7 & 0,245 & Alta & 1,0 & 565 & 0,35 & 198 \\
R-24 & 0,35 & 0,7 & 0,245 & Alta & 0,0 & 480 & 0,00 & 168 \\
Total & — & — & 5,292869 & — & — & — & 10,136475 & 4.746,974 \\
\end{tabla}
La exposición inherente suma unos 5,29 con las probabilidades puntuales del registro, y 5,06 sin R-15, que es el escenario conjunto de otros cinco riesgos; la Tabla \ref{tab:exposicion-valor-esperado-inherentes-24-riesgos} presenta los valores redondeados. R-15 usa una probabilidad de 45,6 \%, calculada con las 32 combinaciones del Anexo D. Los valores esperados de plazo corresponden a tramos diferentes y su suma no representa un atraso único.

\seccion{Exposición residual}
La Tabla \ref{tab:exposicion-valor-esperado-residuales-24-riesgos} presenta los mismos riesgos después de aplicar las respuestas del Formulario T-16, con la probabilidad, el impacto, las semanas y las horas que quedan.

\begin{tabla}{I{0.081}D{0.137}D{0.062}D{0.118}I{0.105}D{0.130}D{0.099}D{0.130}D{0.137}}{Exposición y valor esperado residuales de los 24 riesgos.}{tab:exposicion-valor-esperado-residuales-24-riesgos}
  {\cab{ID} & \cab{P} & \cab{I} & \cab{E = P × I} & \cab{Zona} & \cab{Semanas} & \cab{HH} & \cab{VE plazo (sem)} & \cab{VE esfuerzo (HH)}}
  {Fuente: elaboración propia a partir de Formulario T-16 (asociado al cap. 8), sección 2.}
R-01 & 0,20 & 0,3 & 0,060 & Baja & 0,5 & 300 & 0,10 & 60 \\
R-02 & 0,15 & 0,5 & 0,075 & Baja & 0,0 & 800 & 0,00 & 120 \\
R-03 & 0,15 & 0,5 & 0,075 & Baja & 1,0 & 400 & 0,15 & 60 \\
R-04 & 0,10 & 0,7 & 0,070 & Baja & 0,0 & 300 & 0,00 & 30 \\
R-05 & 0,20 & 0,5 & 0,100 & Media & 1,0 & 200 & 0,20 & 40 \\
R-06 & 0,15 & 0,5 & 0,075 & Baja & 1,0 & 480 & 0,15 & 72 \\
R-07 & 0,20 & 0,5 & 0,100 & Media & 1,0 & 240 & 0,20 & 48 \\
R-08 & 0,30 & 0,5 & 0,150 & Media & 0,0 & 100 & 0,00 & 30 \\
R-09 & 0,15 & 0,5 & 0,075 & Baja & 1,0 & 100 & 0,15 & 15 \\
R-10 & 0,20 & 0,5 & 0,100 & Media & 1,0 & 310 & 0,20 & 62 \\
R-11 & 0,55 & 0,3 & 0,165 & Media & 0,0 & 0 & 0,00 & 0 \\
R-12 & 0,20 & 0,7 & 0,140 & Media & 2,0 & 1.130 & 0,40 & 226 \\
R-13 & 0,15 & 0,3 & 0,045 & Baja & 0,0 & 160 & 0,00 & 24 \\
R-14 & 0,20 & 0,5 & 0,100 & Media & 0,0 & 160 & 0,00 & 32 \\
R-15 & 0,0017600 & 0,3 & 0,000528 & Baja & 0,3 & 100 & 0,000528 & 0,176 \\
R-16 & 0,25 & 0,7 & 0,175 & Media & 0,0 & 400 & 0,00 & 100 \\
R-17 & 0,15 & 0,7 & 0,105 & Media & 0,0 & 384 & 0,00 & 58 \\
R-18 & 0,10 & 0,5 & 0,050 & Baja & 1,0 & 240 & 0,10 & 24 \\
R-19 & 0,10 & 0,7 & 0,070 & Baja & 2,0 & 684 & 0,20 & 68 \\
R-20 & 0,10 & 0,7 & 0,070 & Baja & 0,0 & 400 & 0,00 & 40 \\
R-21 & 0,20 & 0,3 & 0,060 & Baja & 0,0 & 200 & 0,00 & 40 \\
R-22 & 0,10 & 0,7 & 0,070 & Baja & 0,0 & 160 & 0,00 & 16 \\
R-23 & 0,10 & 0,3 & 0,030 & Baja & 0,5 & 280 & 0,05 & 28 \\
R-24 & 0,10 & 0,5 & 0,050 & Baja & 0,0 & 480 & 0,00 & 48 \\
Total & — & — & 2,010528 & — & — & — & 1,900528 & 1.241,176 \\
\end{tabla}
La exposición agregada baja de 5,06 a 2,01 sin contar R-15, para no sumar dos veces sus cinco riesgos: una reducción estimada de 60 \%. Quince riesgos quedan en zona baja y nueve en media. R-16 tiene la mayor exposición residual, 0,175. La suma de valores esperados de esfuerzo por ficha es 1.241,176 HH; no es la reserva de contingencia, que excluye las aceptaciones pasivas, R-15 y considera O-02. Los cálculos emplean los valores completos, aunque las tablas presenten redondeos.

\seccion{Síntesis por zona y por rama}
La Tabla \ref{tab:distribucion-riesgos-zona-inherente-residual} resume cuántos riesgos hay en cada zona antes y después de las respuestas, y la Tabla \ref{tab:exposicion-rama-estructura-desglose-riesgos} cómo se reparte la exposición entre las ramas de la estructura de desglose.

\begin{tabla}{I{0.112}D{0.137}I{0.311}D{0.130}I{0.311}}{Distribución de los riesgos por zona, inherente y residual.}{tab:distribucion-riesgos-zona-inherente-residual}
  {\cab{Zona} & \cab{N° inherente} & \cab{Riesgos (inherente)} & \cab{N° residual} & \cab{Riesgos (residual)}}
  {Fuente: elaboración propia a partir de Formulario T-16 (asociado al cap. 8), sección 2.}
Crítica & 3 & R-05, R-08, R-12 & 0 & Ninguno \\
Alta & 10 & R-01, R-07, R-11, R-15, R-16, R-19, R-20, R-21, R-23, R-24 & 0 & Ninguno \\
Media & 8 & R-04, R-06, R-09, R-10, R-13, R-14, R-17, R-18 & 9 & R-05, R-07, R-08, R-10, R-11, R-12, R-14, R-16, R-17 \\
Baja & 3 & R-02, R-03, R-22 & 15 & R-01, R-02, R-03, R-04, R-06, R-09, R-13, R-15, R-18, R-19, R-20, R-21, R-22, R-23, R-24 \\
\end{tabla}
\begin{tabla}{I{0.348}D{0.143}D{0.180}D{0.143}D{0.186}}{Exposición por rama de la estructura de desglose de riesgos.}{tab:exposicion-rama-estructura-desglose-riesgos}
  {\cab{Rama} & \cab{N° de riesgos} & \cab{Exposición inherente} & \cab{\% del total} & \cab{Exposición residual}}
  {Fuente: elaboración propia a partir de Formulario T-16 (asociado al cap. 8), sección 2.}
1 · Solución y tecnología & 5 & 0,765000 & 14,5 \% & 0,350000 \\
2 · Integración y legado & 2 & 0,525000 & 9,9 \% & 0,175000 \\
3 · Operación del cliente & 4 & 1,085000 & 20,5 \% & 0,425000 \\
4 · Personas y adopción & 4 & 1,020000 & 19,3 \% & 0,450000 \\
5 · Contrato y licitación & 2 & 0,297869 & 5,6 \% & 0,070528 \\
6 · Regulatorio y externo & 4 & 0,935000 & 17,7 \% & 0,410000 \\
7 · Empresa proveedora & 3 & 0,665000 & 12,6 \% & 0,130000 \\
Total & 24 & 5,292869 & 100 \% & 2,010528 \\
\end{tabla}
De las tablas se desprende que la operación del cliente y las personas y adopción reúnen el 40 \% de la exposición inherente, y que todas las ramas bajan más de la mitad con las respuestas.

\seccion{Oportunidades}
Las cuatro oportunidades se miden con las mismas escalas, donde el impacto representa el beneficio. La Tabla \ref{tab:exposicion-oportunidades} las presenta.

\begin{tabla}{I{0.087}I{0.596}D{0.099}D{0.099}D{0.118}}{Exposición de las oportunidades.}{tab:exposicion-oportunidades}
  {\cab{ID} & \cab{Oportunidad} & \cab{P} & \cab{I} & \cab{E = P × I}}
  {Fuente: elaboración propia a partir de Formulario T-16 (asociado al cap. 8), sección 3.}
O-01 & Menos embarques penalizados por arsénico desde 2028 & 0,50 & 0,7 & 0,350 \\
O-02 & Fabricante que entrega la telemetría por ISO 15143-3 & 0,40 & 0,3 & 0,120 \\
O-03 & Replicación de SITCO en la operación que el cliente evalúa adquirir en 2030 & 0,25 & 0,5 & 0,125 \\
O-04 & Recuperar parte de los 47 minutos perdidos por cambio de turno & 0,35 & 0,5 & 0,175 \\
\end{tabla}
De la tabla se desprende que la oportunidad de mayor peso es reducir los embarques penalizados por arsénico (O-01), que depende de la innovación 1 y del acuerdo con el operador portuario. Sólo O-02 tiene un efecto en horas, de 160 HH, cuyo valor esperado de 64 HH se resta de la reserva de contingencia (Anexo E).

\anexo{Análisis de modos y efectos de falla}
Este anexo detalla el análisis que la sección 8.2.4 del Subdocumento 8 resume en la Tabla \refext{S8}{tab:numero-prioridad-riesgos-zona-critica-alta}. Se aplica a los trece riesgos de zona crítica y alta, siguiendo el método de la norma IEC 60812:2018, y sus tres escalas se derivan del registro para no introducir una segunda opinión sobre la probabilidad y el impacto.

\seccion{Escalas de severidad, ocurrencia y detección}
La severidad (S) se toma del nivel de impacto y la ocurrencia (O) de la probabilidad, redondeada a la decena (O = P × 10). La detección (D) es la única escala nueva: mide cuán tarde aparece el disparador respecto del efecto, como muestra la Tabla \ref{tab:correspondencia-severidad-escala-deteccion}.

\begin{tabla}{D{0.095}I{0.234}I{0.671}}{Correspondencia de severidad y escala de detección.}{tab:correspondencia-severidad-escala-deteccion}
  {\cab{Valor} & \cab{Severidad (S) según el impacto} & \cab{Detección (D) según el disparador}}
  {Fuente: elaboración propia a partir de Subdocumento 8, cap. 8, sección 8.2.4.}
2 & Muy bajo (0,1) & Disparador con fecha conocida más de tres meses antes del efecto \\
3 & No se usa & Disparador con fecha fija en los primeros meses, antes de afectar la ruta \\
4 & Bajo (0,3) & Disparador medido cada mes con un indicador cuantitativo \\
5 & No se usa & Disparador que aparece en la puesta en marcha de un componente \\
6 & Medio (0,5) & Disparador una o varias semanas antes del hito o turno afectado \\
7 & No se usa & Disparador que aparece durante la marcha blanca \\
8 & Alto (0,7) & Falla que sólo se detecta en la aceptación del hito \\
10 & Muy alto (0,9) & Sin disparador observable \\
\end{tabla}
De la tabla se desprende que un D alto no significa que el riesgo sea más probable, sino que deja menos tiempo para reaccionar. Por eso el RPN complementa la exposición y no la reemplaza: la priorización contractual sigue siendo la de la matriz P × I.

\seccion{Aplicación a los riesgos de zona crítica y alta}
La Tabla \ref{tab:modos-falla-disparadores-numero-prioridad} presenta el modo de falla de cada riesgo, el disparador que lo detecta y su número de prioridad RPN = S × O × D.

\begin{tabla}{I{0.106}I{0.323}I{0.335}D{0.056}D{0.056}D{0.056}D{0.068}}{Modos de falla, disparadores y número de prioridad.}{tab:modos-falla-disparadores-numero-prioridad}
  {\cab{Riesgo} & \cab{Modo de falla} & \cab{Disparador que lo detecta} & \cab{S} & \cab{O} & \cab{D} & \cab{RPN}}
  {Fuente: elaboración propia a partir de Subdocumento 8, cap. 8, sección 8.2.4.}
R-12 & Los operadores siguen declarando por radio al cierre de la marcha blanca & Declaraciones propias bajo 90 \% en la segunda semana de la pala piloto (marcha blanca) & 8 & 6 & 7 & 336 \\
R-08 & Una pala trabaja más de 48 horas sin señal y se pierden registros & Estudio de cobertura del mes 4 o equipo con más de 24 horas sin sincronizar & 8 & 6 & 4 & 192 \\
R-23 & Personal sin credencial deja turnos descubiertos y prolonga la marcha blanca & Credencial sin resolver a 5 días hábiles del primer turno o examen de altura reprobado & 8 & 4 & 6 & 192 \\
R-16 & Un fabricante no firma el acceso a la telemetría o pide un pago & Acuerdo sin firmar al cierre del mes 9 & 8 & 5 & 4 & 160 \\
R-05 & Una interfaz crítica no existe al cierre del levantamiento & Documentación faltante al cierre del mes 2 o prueba de lectura fallida en el mes 3 & 8 & 6 & 3 & 144 \\
R-11 & No hay acuerdo sindical para la verificación facial & Sin acuerdo firmado al cierre del mes 12, tras 11 sesiones mensuales desde M2; las 13 sesiones previstas se completan en M14 & 6 & 6 & 4 & 144 \\
R-19 & La primera campaña de intrusión deja hallazgos altos abiertos & Revisión interna del mes 9, dos meses antes de la campaña & 8 & 3 & 6 & 144 \\
R-07 & Los factores de ajuste quedan sin explicación o la fuente no se ratifica & Acta del mes 3 sin firma de finanzas o factor sin explicar al cierre del mes 2 & 8 & 5 & 3 & 120 \\
R-21 & La segunda fibra no opera o el satélite rinde menos de 20 Mbps & Fecha o avance incompatible con M6, ruta compartida en revisión previa o satélite inferior a 20 Mbps en M5 & 6 & 4 & 5 & 120 \\
R-24 & Faltan analistas al inicio de la operación 24x7 & Conteo de analistas al cierre de los meses 12 y 18 & 8 & 4 & 3 & 96 \\
R-15 & La cadena crítica de la Etapa 1 supera su reserva de plazo & Consumo de la reserva revisado en cada Comité de Proyecto & 6 & 5 & 3 & 90 \\
R-01 & Un componente pierde soporte en plena construcción & Fechas públicas de fin de soporte, con aviso a 90 días & 6 & 5 & 2 & 60 \\
R-20 & El cliente no autoriza la réplica de datos sensibles fuera de Chile & Autorización pedida con H1 y exigida al cierre del mes 4 & 8 & 3 & 2 & 48 \\
\end{tabla}
Cuatro riesgos superan RPN 150: R-12 (336), R-08 y R-23 (192), y R-16 (160). En R-23, el disparador a cinco días hábiles del primer turno equivale aproximadamente a una semana laboral de anticipación: D=6 conforme a la Tabla \ref{tab:correspondencia-severidad-escala-deteccion}. La fecha de ingreso es conocida, pero la credencial sin resolver no se detecta con meses de margen; un examen reprobado activa antes el mismo control. Para responder se mantienen personas de reserva acreditadas. La detección temprana también incluye la prueba con guantes de M5 para R-12, el estudio de cobertura para R-08 y la notificación a fabricantes en el Comité Ejecutivo de M1 para R-16.

\anexo{Simulación de la cadena crítica de la Etapa 1}
Este anexo detalla la simulación que la sección 8.2.5 del Subdocumento 8 resume en la Figura \refext{S8}{fig:distribucion-atraso-acumulado-sobre-cadena-critica} y en la Tabla \refext{S8}{tab:resultado-simulacion-sobre-cadena-critica-etapa}. Su objetivo es estimar la probabilidad de que los riesgos que caen sobre la cadena crítica de la Etapa 1 consuman la reserva de plazo del mes 11, de 4,3 semanas, declarada en el paquete 1.1.3.2 del Formulario T-14.

\seccion{Modelo y parámetros}
Cada riesgo de la cadena se modela como un evento de Bernoulli: ocurre con su probabilidad y, si ocurre, suma sus semanas brutas de atraso. El atraso total es la suma de los eventos que ocurren. R-15 no se suma como evento, porque representa precisamente el escenario conjunto de los otros cinco. Los cinco eventos caen sobre la cadena A04 a A24 de la Tabla \refext{T14}{tab:red-precedencias-calculo-ruta-critica-meses} del Formulario T-14, que no tiene holgura. R-07 llega a ella porque la Conciliación Metalúrgica (1.4.3.4) está programada en los meses 6 a 11, sus fechas más tardías, y su retrabajo pasa a la corrección del mes 11 (A16). La Tabla \ref{tab:parametros-simulacion} presenta los parámetros.

\begin{tabla}{I{0.441}D{0.130}D{0.149}D{0.130}D{0.149}}{Parámetros de la simulación.}{tab:parametros-simulacion}
  {\cab{Riesgo} & \cab{P inherente} & \cab{Semanas inherentes} & \cab{P residual} & \cab{Semanas residuales}}
  {Fuente: elaboración propia a partir de Formulario T-16 (asociado al cap. 8), sección 2; Formulario T-14 (asociado al cap. 7), tabla \refext{T14}{tab:red-precedencias-calculo-ruta-critica-meses} del Formulario T-14 (red de precedencias).}
R-01 · Fin de soporte de componentes en la Etapa 1 & 0,45 & 1,5 & 0,20 & 0,5 \\
R-05 · Interfaz existente no documentada & 0,60 & 2,5 & 0,20 & 1,0 \\
R-07 · Factores de ajuste sin quien los explique & 0,45 & 1,5 & 0,20 & 1,0 \\
R-18 · Salida o sobrecarga de una persona clave & 0,25 & 2,0 & 0,10 & 1,0 \\
R-19 · Hallazgos altos en las pruebas de intrusión & 0,30 & 3,0 & 0,10 & 2,0 \\
\end{tabla}
De la tabla se desprende que R-05 es el evento de mayor peso, con 60 \% de probabilidad y 2,5 semanas. Se ejecutaron 100.000 iteraciones con la semilla fija 20261010, para que el resultado sea reproducible, y se calculó además la distribución exacta de las 32 combinaciones posibles de los cinco eventos. La diferencia entre ambos métodos es menor que 0,1 punto porcentual en todos los indicadores.

\seccion{Resultados}
La Tabla \ref{tab:distribucion-exacta-atraso-sobre-cadena-critica} presenta la distribución exacta del atraso, con su probabilidad acumulada, sin respuestas y con las respuestas del registro.

\begin{tabla}{D{0.162}D{0.211}D{0.211}D{0.211}D{0.205}}{Distribución exacta del atraso sobre la cadena crítica de la Etapa 1.}{tab:distribucion-exacta-atraso-sobre-cadena-critica}
  {\cab{Atraso (sem)} & \cab{Probabilidad sin respuestas} & \cab{Acumulada sin respuestas} & \cab{Probabilidad con respuestas} & \cab{Acumulada con respuestas}}
  {Fuente: elaboración propia.}
0,0 & 6,4 \% & 6,4 \% & 41,5 \% & 41,5 \% \\
0,5 & 0,0 \% & 6,4 \% & 10,4 \% & 51,8 \% \\
1,0 & 0,0 \% & 6,4 \% & 25,3 \% & 77,2 \% \\
1,5 & 10,4 \% & 16,7 \% & 6,3 \% & 83,5 \% \\
2,0 & 2,1 \% & 18,9 \% & 9,5 \% & 93,0 \% \\
2,5 & 9,5 \% & 28,4 \% & 2,4 \% & 95,4 \% \\
3,0 & 7,0 \% & 35,4 \% & 3,1 \% & 98,5 \% \\
3,5 & 3,5 \% & 38,8 \% & 0,8 \% & 99,3 \% \\
4,0 & 15,6 \% & 54,4 \% & 0,5 \% & 99,8 \% \\
4,5 & 7,6 \% & 62,1 \% & 0,1 \% & 100,0 \% \\
5,0 & 2,3 \% & 64,4 \% & 0,0 \% & 100,0 \% \\
5,5 & 10,5 \% & 74,8 \% & 0,0 \% & 100,0 \% \\
6,0 & 7,0 \% & 81,9 \% & 0,0 \% & 100,0 \% \\
6,5 & 1,5 \% & 83,4 \% & 0,0 \% & 100,0 \% \\
7,0 & 6,7 \% & 90,0 \% & 0,0 \% & 100,0 \% \\
7,5 & 3,5 \% & 93,5 \% & 0,0 \% & 100,0 \% \\
8,0 & 0,6 \% & 94,1 \% & 0,0 \% & 100,0 \% \\
8,5 & 2,7 \% & 96,9 \% & 0,0 \% & 100,0 \% \\
9,0 & 2,2 \% & 99,1 \% & 0,0 \% & 100,0 \% \\
10,5 & 0,9 \% & 100,0 \% & 0,0 \% & 100,0 \% \\
\end{tabla}
De la tabla se desprende que, sin respuestas, el atraso acumulado supera el 80 \% de los casos recién en 6,0 semanas; con las respuestas, el 80 \% de los casos queda en 1,5 semanas o menos. La Tabla \ref{tab:probabilidad-superar-umbrales-reserva-plazo} compara la probabilidad de superar los umbrales de escalamiento 3 y 4 del Anexo A y la reserva completa.

\begin{tabla}{D{0.149}I{0.453}D{0.199}D{0.199}}{Probabilidad de superar los umbrales de la reserva de plazo.}{tab:probabilidad-superar-umbrales-reserva-plazo}
  {\cab{Atraso} & \cab{Significado} & \cab{Sin respuestas} & \cab{Con respuestas}}
  {Fuente: elaboración propia.}
2,2 sem & Umbral 3: 50 \% de la reserva & 81,1 \% & 7,0 \% \\
3,4 sem & Umbral 4: 80 \% de la reserva & 64,6 \% & 1,5 \% \\
4,3 sem & Reserva completa del mes 11 & 45,6 \% & 0,2 \% \\
\end{tabla}
De la tabla se desprende que, con las respuestas, la probabilidad de superar la reserva completa es de 0,2 \%, muy por debajo del umbral 5 de 20 \%, y que la reserva de contingencia de plazo, con valor esperado de 0,8 semanas, deja 2,8 semanas al P80 como reserva de gestión de plazo.

\seccion{Alcance y límites del modelo}
El modelo considera los cinco eventos identificados que afectan la cadena crítica de la Etapa 1. Supone que ocurren de forma independiente; si comparten causas, sus probabilidades conjuntas deben revisarse antes de usar el resultado para decidir una reserva. No simula la variabilidad de duración de la construcción. También supone que A16 puede absorber atraso comprimiendo corrección, que continúa con 1.4.3.7 en M12. La capacidad utilizable se comprueba con la carga por rol y las condiciones de entrada a certificación. Si sólo una parte de las 4,3 semanas puede comprimirse, la distribución de la tabla se compara con ese margen menor. Si no existe absorción, el atraso alcanza H5. La ventana de siete días protege H7 y no H5. En la Etapa 2, el cuarto puesto de escaneo es la medida prevista frente a baja productividad de B06, con 18,9 semanas esperadas.

\anexo{Reservas de contingencia y de gestión}
Este anexo detalla las reservas que la sección 8.3.3 del Subdocumento 8 resume en la Tabla \refext{S8}{tab:reservas-contingencia-gestion}. La reserva de contingencia es el valor esperado del riesgo residual de cada riesgo con respuesta o aceptado en forma activa, medido en semanas y en horas hombre, y se suma a las 988 HH de incertidumbre de la estimación que el Formulario T-14 cubre con los paquetes de corrección y el saldo de los roles clave. La reserva de gestión cubre los riesgos no identificados. La valorización de ambas va en la Oferta Económica, conforme al Art. 50.2.

\seccion{Reserva de contingencia por riesgo}
La Tabla \ref{tab:reserva-contingencia-riesgo} presenta la reserva de cada riesgo, el tramo del cronograma donde se refleja y el mes en que se libera si el riesgo se cierra sin ocurrir. R-02 y R-03 no tienen reserva porque se aceptan de forma pasiva, y R-15 no se suma porque su reserva ya está en los cinco riesgos que lo componen.

O-02 aporta un ahorro esperado de 64 HH, no un recurso asegurado. El registro conserva la contingencia neta de 997 HH y permite identificar una exposición bruta de 1.061 HH antes de ese descuento. Hasta que la oportunidad se materialice, el Jefe de Proyecto no utiliza su ahorro como capacidad liberada. Si no ocurre, la cobertura se decide con la reserva bruta y con el saldo real de los roles, sin alterar la probabilidad ni el impacto de las amenazas.

\begin{tabla}{I{0.087}I{0.323}D{0.155}D{0.161}I{0.273}}{Reserva de contingencia por riesgo.}{tab:reserva-contingencia-riesgo}
  {\cab{ID} & \cab{Tramo donde se refleja} & \cab{VE plazo (sem)} & \cab{VE esfuerzo (HH)} & \cab{Se libera en}}
  {Fuente: elaboración propia a partir de Formulario T-16 (asociado al cap. 8), sección 2.}
R-01 & Etapa 1 (H5) & 0,10 & 60 & M8 (versiones actualizadas) \\
R-04 & Operación & 0,00 & 30 & Cada campaña anual de intrusión sin hallazgos \\
R-05 & Etapa 1 (H5) & 0,20 & 40 & M4 (interfaces probadas) \\
R-06 & Datos de la Etapa 2 (H10 y H12) & 0,15 & 72 & M15 (último lote) \\
R-07 & Etapa 1 (H5) & 0,20 & 48 & M15 (cierre de marcha blanca E1) \\
R-08 & Implantación E1 (H7) & 0,00 & 30 & M16 (paso a producción E1) \\
R-09 & Sala técnica y H6 & 0,15 & 15 & M13 (H6) \\
R-10 & Implantación E1 y E2 (H7, H12) & 0,20 & 62 & M21 (H12) \\
R-11 & Etapa 2 (H9) & 0,00 & 0 & M15 (acuerdo o plan B) \\
R-12 & Implantación E1 (H7) & 0,40 & 226 & M16 (H7) \\
R-13 & Etapa 2 (H9 a H12) & 0,00 & 24 & M21 (H12) \\
R-14 & Etapa 2 y operación & 0,00 & 32 & M30 (segundo informe de verificación) \\
R-16 & Etapa 2 (H9) & 0,00 & 100 & M13 (acuerdos firmados) \\
R-17 & Operación (M25) & 0,00 & 58 & M25 (enero de 2029) \\
R-18 & Etapa 1 (H5) & 0,10 & 24 & M21 (H12) \\
R-19 & Etapa 1 (H5) & 0,20 & 68 & M20 (última campaña) \\
R-20 & Diseño y continuidad (H2, H7) & 0,00 & 40 & M4 (autorización recibida) \\
R-21 & Implantación E1 (H7) & 0,00 & 40 & M16 (segunda fibra operativa) \\
R-22 & Calendario (H1) & 0,00 & 16 & M1 (contrato firmado) \\
R-23 & Implantación E1 (H7) & 0,05 & 28 & M20 (fin de las marchas blancas) \\
R-24 & Operación & 0,00 & 48 & M24 (tres meses de operación medidos) \\
O-02 & Etapa 2 (H9) & 0,00 & -64 & M17 (H9) \\
Total & Tramos paralelos: el plazo no se suma & 0,80 en la Etapa 1 & 997 & No aplica \\
\end{tabla}
La contingencia neta es de 997 HH después de descontar las 64 HH esperadas de O-02; la exposición bruta es de 1.061 HH. R-04 aporta 30 HH, R-17 58 HH y R-24 48 HH de efectos que alcanzan la operación. R-14 aporta 32 HH con exposición en implementación y operación. La tabla identifica el efecto de cada ficha, pero no reemplaza su asignación mensual ni demuestra que una reserva de implantación cubra las recurrencias posteriores de un riesgo.

\seccion{Reserva por tramo del cronograma}
La Tabla \ref{tab:reserva-contingencia-tramo-uso-reserva-plazo} agrupa la reserva por tramo y la compara con la reserva de plazo que ese tramo ya tiene en el Formulario T-14.

\begin{tabla}{I{0.224}I{0.273}D{0.118}D{0.087}I{0.298}}{Reserva de contingencia por tramo y uso de la reserva de plazo.}{tab:reserva-contingencia-tramo-uso-reserva-plazo}
  {\cab{Tramo} & \cab{Riesgos} & \cab{Plazo (sem)} & \cab{HH} & \cab{Uso de la reserva de plazo}}
  {Fuente: elaboración propia a partir de Formulario T-14 (asociado al cap. 7), tabla \refext{T14}{tab:red-precedencias-calculo-ruta-critica-meses} del Formulario T-14 (red de precedencias).}
Etapa 1 hacia H5 & R-01, R-05, R-07, R-18, R-19 & 0,80 & 240 & 19 \% de las 4,3 sem del mes 11 en valor esperado \\
Implantación E1 (H7) & R-08, R-12, R-21, R-23 & 0,45 & 324 & Se cubre con acompañamiento en faena sin mover fases; 7 días de ventana \\
Implantación E1 y E2 & R-10 & 0,20 & 62 & 7 días de ventana en E1 y 2 días en E2 \\
Datos de la Etapa 2 & R-06 & 0,15 & 72 & 2 días de ventana; palanca: cuarto puesto de escaneo \\
Sala técnica y H6 & R-09 & 0,15 & 15 & Fuera de la cadena crítica \\
Etapa 2 y operación (R-17 en M25) & R-11, R-13, R-16, R-17 & 0 & 182 & Sin efecto en hitos \\
Diseño, calendario y operación & R-04, R-14, R-20, R-22, R-24 & 0 & 166 & Sin efecto en hitos \\
Oportunidad & O-02 & 0 & -64 & Se resta de la reserva \\
Total & 21 riesgos y 1 oportunidad & No se suma & 997 & Asignación por riesgo, persona y mes \\
\end{tabla}
La Tabla \ref{tab:reserva-contingencia-tramo-uso-reserva-plazo} distribuye las 997 HH netas por grupos de riesgos, sin certificar una separación mensual de capacidad. La autorización identifica riesgo, persona, rol, mes, paquete ejecutor y saldo, y distingue trabajo preventivo de contingencia. Los paquetes 1.4.3.7 y 1.4.4.7 suman 1.280 HH y mantienen su corrección y cobertura de incertidumbre del T-14. La misma hora no cubre simultáneamente defectos, incertidumbre y riesgo. Las 0,8 semanas son atraso medio; la absorción en M11 depende de la corrección que pueda comprimirse.

\seccion{Reserva de gestión}
La reserva de gestión cubre lo que el registro no identificó y sólo la libera el Comité Ejecutivo. La Tabla \ref{tab:componentes-reserva-gestion} presenta sus tres componentes.

\begin{tabla}{I{0.199}I{0.199}I{0.373}I{0.230}}{Componentes de la reserva de gestión.}{tab:componentes-reserva-gestion}
  {\cab{Componente} & \cab{Monto} & \cab{Criterio} & \cab{Liberación}}
  {Fuente: elaboración propia a partir de Formulario T-14 (asociado al cap. 7), tabla \refext{T14}{tab:red-precedencias-calculo-ruta-critica-meses} del Formulario T-14 (red de precedencias).}
Plazo, Etapa 1 & 2,8 semanas & Reserva del mes 11 (4,3 sem) menos el P80 residual (1,5 sem) & Con la aceptación de H5 \\
Plazo, ventanas de paso & 7 días en E1 y 2 días en E2 & Margen después de la última ola prevista (hasta el 28 de abril y el 29 y 30 de septiembre de 2028) & Con H7 y con H12 \\
Esfuerzo & 5.372 HH & Desviación de la estimación si todas sus desviaciones fueran en el mismo sentido (Formulario T-14), 8,8 \% de las 61.194 HH de las ramas 1.1 a 1.7; 3.223 HH para la Etapa 1 y 2.149 HH para la Etapa 2 & Con las actas de cierre de cada etapa (meses 16 y 21) \\
\end{tabla}
De la tabla se desprende que el margen calculado de 2,8 semanas resulta de restar el P80 residual de 1,5 semanas a las 4,3 semanas consideradas para M11. Su uso depende de que esa capacidad de absorción sea demostrable. Sin respuestas, el P80 de 6,0 semanas supera el periodo considerado y no deja ese margen. Las ventanas de paso protegen únicamente sus propios hitos y no se trasladan a una aceptación anterior.

\seccion{Política de uso y liberación}
El uso de las reservas sigue cuatro reglas, que el Jefe de Proyecto informa en cada Comité de Proyecto:

La reserva de contingencia la autoriza el Jefe de Proyecto durante la implementación y el Gerente de Servicio durante la operación desde M21, contra un riesgo cuyo disparador se activó. Se registra causa, periodo, semanas, horas, rol ejecutor y saldo utilizado. La reserva de gestión la libera el Comité Ejecutivo.

Cada línea se libera cuando se verifica el cierre de su riesgo. Las fechas del registro son fechas previstas de revisión, no liberaciones automáticas. El saldo no se reasigna sin aprobación del Comité de Proyecto.

El consumo de 50 \% de una reserva obliga a un plan de recuperación, y el de 80 \% la escala al Comité Ejecutivo.

El uso de una reserva no modifica por sí mismo ningún hito contractual. Se comprueba el margen utilizable y se registra la asignación de horas por rol y paquete, sin superponerla con corrección o incertidumbre de estimación. Si no permite cumplir el hito, la desviación se escala y se gestiona conforme al procedimiento contractual.

\anexo{Análisis costo-beneficio de las respuestas}
Este anexo compara el esfuerzo de las respuestas con el cambio del impacto esperado de cada ficha. Mantiene las probabilidades y consecuencias del registro. Trece respuestas se apoyan en paquetes existentes y once requieren 704 HH adicionales respecto de sus actividades de referencia. La ejecución corresponde al rol y al paquete de cada ficha. El plan de respuestas 1.1.3.2 coordina su asignación, pero no demuestra por sí solo que exista capacidad libre. Las medidas preventivas se separan de la contingencia residual y se valorizan en la Oferta Económica. La tabla conserva el detalle por riesgo para permitir su revisión.

\begin{tabla}{I{0.081}I{0.249}D{0.137}D{0.130}D{0.130}D{0.124}D{0.149}}{Costo y beneficio de la respuesta de cada riesgo.}{tab:costo-beneficio-respuesta-cada-riesgo}
  {\cab{ID} & \cab{Estrategia} & \cab{Costo adicional (HH)} & \cab{VE inherente (HH)} & \cab{VE residual (HH)} & \cab{Red. HH} & \cab{Beneficio / costo}}
  {Fuente: elaboración propia a partir de Formulario T-16 (asociado al cap. 8), sección 2; Formulario T-14 (asociado al cap. 7), tabla \refext{T14}{tab:esfuerzo-estimado-rama-origen-horas} del Formulario T-14 (esfuerzo estimado por rama).}
R-01 & Mitigar & 240 & 225 & 60 & 165 & 0,7 \\
R-02 & Aceptar & 0 & 120 & 120 & 0 & Sin costo adicional \\
R-03 & Aceptar & 0 & 60 & 60 & 0 & Sin costo adicional \\
R-04 & Mitigar & 0 & 60 & 30 & 30 & Sin costo adicional \\
R-05 & Mitigar & 64 & 240 & 40 & 200 & 3,1 \\
R-06 & Aceptar & 0 & 72 & 72 & 0 & Sin costo adicional \\
R-07 & Escalar & 16 & 180 & 48 & 132 & 8,3 \\
R-08 & Mitigar y escalar la inversión en red al cliente & 0 & 678 & 30 & 648 & Sin costo adicional \\
R-09 & Transferir & 40 & 50 & 15 & 35 & 0,9 \\
R-10 & Escalar & 16 & 109 & 62 & 47 & 2,9 \\
R-11 & Escalar & 0 & 211 & 0 & 211 & Sin costo adicional \\
R-12 & Mitigar & 0 & 1.017 & 226 & 791 & Sin costo adicional \\
R-13 & Mitigar & 16 & 48 & 24 & 24 & 1,5 \\
R-14 & Mitigar & 0 & 56 & 32 & 24 & Sin costo adicional \\
R-15 & Mitigar & 0 & 311,724 & 0,176 & 311,548 & Sin costo adicional \\
R-16 & Escalar & 0 & 200 & 100 & 100 & Sin costo adicional \\
R-17 & Escalar & 40 & 96 & 58 & 38 & 1,0 \\
R-18 & Mitigar & 0 & 60 & 24 & 36 & Sin costo adicional \\
R-19 & Mitigar & 40 & 308 & 68 & 239 & 6,0 \\
R-20 & Escalar & 40 & 120 & 40 & 80 & 2,0 \\
R-21 & Mitigar & 0 & 80 & 40 & 40 & Sin costo adicional \\
R-22 & Escalar y transferir & 0 & 80 & 16 & 64 & Sin costo adicional \\
R-23 & Mitigar & 32 & 198 & 28 & 170 & 5,3 \\
R-24 & Mitigar & 160 & 168 & 48 & 120 & 0,8 \\
Total & No aplica & 704 & 4.746,974 & 1.241,176 & 3.505,798 & 5,0 \\
\end{tabla}
Las fichas suman 4.746,97445 HH inherentes y 1.241,176 HH residuales, diferencia de 3.505,79845 HH. El cociente descriptivo sobre 704 HH es 5,0; no se interpreta como ahorro independiente porque R-15 representa un escenario conjunto. Su reducción de 311,54845 HH se excluye de cualquier suma de beneficio económico conjunto. Las once medidas con costo adicional conservan 704 HH, reducción de 1.250 HH y cociente 1,8. Se agregan por separado 96 HH para el control transversal de continuidad del Anexo G.4, sin atribuirle una reducción probabilística no medida.\newline El costo preventivo adicional de R-12 es 0 HH porque sus acciones se ejecutan en paquetes ya incluidos en la EDT. Las HH de G.3 corresponden al esfuerzo contingente para extender la marcha blanca si se activa el riesgo; no forman parte del costo preventivo comparado en esta tabla.

\begin{horizontal}

\anexo{Decisiones y recursos de las respuestas}
Este anexo permite reproducir las valoraciones y asignar los recursos adicionales. Los planes son estimaciones de ejecución vinculadas al T-14, el T-15 y los requisitos de aceptación; la contratación de apoyos y los resultados de los ensayos se registran en sus paquetes antes de habilitar cada servicio. Se separan los valores esperados del registro, el esfuerzo del evento activado y la capacidad de los ejecutores.

\seccion{Aplicación de la escala de impacto}
La Tabla \ref{tab:dimensiones-valoraciones-adoptadas} aplica I = máximo de plazo, HH y alcance. Los días se obtienen de las semanas del registro por siete. Los umbrales exactos son 2, 7, 20 y 30 días, y 300, 600, 1.200 y 3.000 HH. Las consecuencias de R-15 se conservan redondeadas a 2,0 y 0,3 semanas; sus probabilidades se calculan sin redondear. Una semana completa exige al menos impacto medio. El impacto alto de R-12 responde a la extensión fuera del cierre M15, y los de R-16/R-17/R-20/R-22 al requisito esencial aún afectado. R-23 también es alto antes de las respuestas por su semana adicional de marcha blanca; su escenario residual de media semana con personal de reserva mantiene impacto bajo.

\begin{tabla}{I{0.088}I{0.111}I{0.111}I{0.124}I{0.111}I{0.111}I{0.111}I{0.124}I{0.111}}{Dimensiones y valoraciones adoptadas.}{tab:dimensiones-valoraciones-adoptadas}
  {\cab{Riesgo} & \cab{Días inh.} & \cab{HH inh.} & \cab{I alcance inh.} & \cab{I final inh.} & \cab{Días res.} & \cab{HH res.} & \cab{I alcance res.} & \cab{I final res.}}
  {Fuente: elaboración propia a partir de Formulario T-16 (asociado al cap. 8), sección 2.}
R-01 & 10,5 & 500 & 0,3 & 0,5 & 3,5 & 300 & 0,3 & 0,3 \\
R-02 & 0,0 & 800 & 0,5 & 0,5 & 0,0 & 800 & 0,5 & 0,5 \\
R-03 & 7,0 & 400 & 0,5 & 0,5 & 7,0 & 400 & 0,5 & 0,5 \\
R-04 & 0,0 & 300 & 0,9 & 0,9 & 0,0 & 300 & 0,7 & 0,7 \\
R-05 & 17,5 & 400 & 0,7 & 0,7 & 7,0 & 200 & 0,5 & 0,5 \\
R-06 & 21,0 & 480 & 0,3 & 0,7 & 7,0 & 480 & 0,3 & 0,5 \\
R-07 & 10,5 & 400 & 0,7 & 0,7 & 7,0 & 240 & 0,5 & 0,5 \\
R-08 & 14,0 & 1.130 & 0,7 & 0,7 & 0,0 & 100 & 0,5 & 0,5 \\
R-09 & 14,0 & 200 & 0,7 & 0,7 & 7,0 & 100 & 0,5 & 0,5 \\
R-10 & 10,5 & 310 & 0,5 & 0,5 & 7,0 & 310 & 0,5 & 0,5 \\
R-11 & 0,0 & 384 & 0,5 & 0,5 & 0,0 & 0 & 0,3 & 0,3 \\
R-12 & 21,0 & 1.695 & 0,7 & 0,7 & 14,0 & 1.130 & 0,7 & 0,7 \\
R-13 & 0,0 & 160 & 0,5 & 0,5 & 0,0 & 160 & 0,3 & 0,3 \\
R-14 & 0,0 & 160 & 0,5 & 0,5 & 0,0 & 160 & 0,5 & 0,5 \\
R-15 & 14,0 & 684 & 0,3 & 0,5 & 2,1 & 100 & 0,3 & 0,3 \\
R-16 & 0,0 & 400 & 0,7 & 0,7 & 0,0 & 400 & 0,7 & 0,7 \\
R-17 & 0,0 & 384 & 0,7 & 0,7 & 0,0 & 384 & 0,7 & 0,7 \\
R-18 & 14,0 & 240 & 0,7 & 0,7 & 7,0 & 240 & 0,5 & 0,5 \\
R-19 & 21,0 & 1.026 & 0,7 & 0,7 & 14,0 & 684 & 0,7 & 0,7 \\
R-20 & 0,0 & 400 & 0,7 & 0,7 & 0,0 & 400 & 0,7 & 0,7 \\
R-21 & 0,0 & 200 & 0,5 & 0,5 & 0,0 & 200 & 0,3 & 0,3 \\
R-22 & 0,0 & 800 & 0,7 & 0,7 & 0,0 & 160 & 0,7 & 0,7 \\
R-23 & 7,0 & 565 & 0,7 & 0,7 & 3,5 & 280 & 0,3 & 0,3 \\
R-24 & 0,0 & 480 & 0,7 & 0,7 & 0,0 & 480 & 0,5 & 0,5 \\
\end{tabla}
Los valores ordinales y el umbral metodológico de siete días son una decisión de Overfore; no una exigencia numérica de las Bases.

R-05/R-06 conservan una semana residual y pasan a I = 0,5. R-12 conserva dos semanas de extensión de marcha blanca y pasa a I = 0,7. R-16 mantiene I = 0,7 mientras la descarga diaria o los datos de prueba no acrediten telemetría productiva con la frescura exigida. R-17 conserva I = 0,7: SUP-21 permite identificar transitoriamente un origen anterior, pero exige extinguir esa salvedad antes de enero de 2029. R-20 mantiene I = 0,7 hasta demostrar la alternativa de continuidad. R-22 mantiene I = 0,7 porque reprogramar hitos relativos no cambia la fecha comercial absoluta. Estos criterios no sustituyen la aceptación del cliente o del comprador.

R-15 usa P inherente 0,4557375 y P residual 0,00176, suma exacta de las combinaciones que exceden 4,3 semanas. Su exposición residual es 0,00176 × 0,3 = 0,000528. Para lectura se muestra 0,2 \%, pero todos los productos emplean 0,176 \%. El exceso condicionado medio, 1,973965 y 0,336364 semanas, permite comprobar el orden de magnitud de las consecuencias redondeadas. R-15 no es un sexto evento ni una reserva adicional. Las otras probabilidades son estimaciones de planificación del registro y no frecuencias históricas medidas.

\seccion{Programación de las medidas preventivas}
Las once medidas adicionales suman 704 HH y se programan con apoyos nuevos o incrementos de dedicación parcial. No se asignan como horas libres del JP, LD o LF, cuya capacidad base ya se completa con trabajo continuo. La Tabla \ref{tab:asignacion-704-hh-preventivas} identifica el incremento por rol y periodo; el archivo de control por persona y mes conserva la reconstrucción uniforme de la base y sus límites.

\begin{tabla}{I{0.085}I{0.085}I{0.128}I{0.085}I{0.617}}{Asignación de las 704 HH preventivas.}{tab:asignacion-704-hh-preventivas}
  {\cab{Riesgo} & \cab{Mes} & \cab{Rol ejecutor} & \cab{HH} & \cab{Forma de cobertura}}
  {Fuente: elaboración propia.}
R-05 & M2 & QA & 32 & QA-P01; apoyo adicional \\
R-07 & M2 & LDA & 16 & Benjamín Velásquez Farías \\
R-09 & M2 & AS & 16 & Benjamín Rodríguez Pavez \\
R-10 & M2 & ACP & 16 & ACP-P01; apoyo adicional \\
R-05 & M3 & QA & 32 & QA-P01; apoyo adicional \\
R-17 & M3 & ACP & 16 & ACP-P01; apoyo adicional \\
R-17 & M3 & ABG & 24 & ABG-P01; apoyo adicional \\
R-20 & M4 & AS & 24 & Benjamín Rodríguez Pavez \\
R-20 & M4 & BCP & 16 & BCP-P01; apoyo adicional \\
R-09 & M5 & DEV & 24 & DEV-P01; apoyo adicional \\
R-01 & M6 & DEV & 16 & DEV-P01; apoyo adicional \\
R-01 & M7 & DEV & 16 & DEV-P01; apoyo adicional \\
R-01 & M8 & DEV & 16 & DEV-P01; apoyo adicional \\
R-01 & M9 & DEV & 16 & DEV-P01; apoyo adicional \\
R-19 & M9 & ESI & 40 & David Martínez Canto \\
R-01 & M10 & DEV & 16 & DEV-P01; apoyo adicional \\
R-01 & M11 & DEV & 16 & DEV-P01; apoyo adicional \\
R-01 & M12 & DEV & 16 & DEV-P01; apoyo adicional \\
R-23 & M12 & ACP & 32 & ACP-P01; apoyo adicional \\
R-01 & M13 & DEV & 16 & DEV-P01; apoyo adicional \\
R-13 & M13 & RTC & 16 & RTC-P01; apoyo adicional \\
R-01 & M14 & DEV & 16 & DEV-P01; apoyo adicional \\
R-01 & M15 & DEV & 16 & DEV-P01; apoyo adicional \\
R-01 & M16 & DEV & 16 & DEV-P01; apoyo adicional \\
R-24 & M16 & OPS & 32 & OPS-P01; apoyo adicional \\
R-01 & M17 & DEV & 16 & DEV-P01; apoyo adicional \\
R-24 & M17 & OPS & 24 & OPS-P02; apoyo adicional \\
R-01 & M18 & DEV & 16 & DEV-P01; apoyo adicional \\
R-01 & M19 & DEV & 16 & DEV-P01; apoyo adicional \\
R-24 & M19 & OPS & 42 & OPS-P03; apoyo adicional \\
R-01 & M20 & DEV & 16 & DEV-P01; apoyo adicional \\
R-24 & M20 & OPS & 62 & OPS-P04; apoyo adicional \\
\end{tabla}
La asignación nominada sólo incrementa perfiles con margen estimado: Datos en M2, Arquitectura en M2/M4 y Seguridad en M9. Las actividades adicionales de desarrollo, pruebas, coordinación, capacitación y especialidad se cubren con perfiles de apoyo contratados para los periodos del plan. Los identificadores de apoyos describen puestos a incorporar, no personas cuya disponibilidad ya esté certificada.

\seccion{Cobertura de la extensión de marcha blanca}
R-12 activa acompañamiento en las áreas que no cierren sus indicadores; las ya aceptadas no repiten la extensión. El escenario principal es un área de mina durante 21 días. Se asignan seis apoyos en tres parejas semanales: cada apoyo cubre siete turnos de 12 horas, 84 HH, y cada pareja cubre día y noche de una semana. La Tabla \ref{tab:esfuerzo-tres-semanas-area} suma 760 HH brutas y separa 12 HH de LF reordenadas dentro de su dedicación: el incremento es 748 HH. La valoración del registro, 1.695 HH, conserva la equivalencia media 3 × 565 del T-14 y cubre este escenario sin certificar que cubra cualquier número de áreas.

\begin{tabla}{I{0.216}I{0.294}I{0.314}I{0.176}}{Esfuerzo de tres semanas para un área.}{tab:esfuerzo-tres-semanas-area}
  {\cab{Trabajo} & \cab{Cálculo} & \cab{Ejecutor} & \cab{HH}}
  {Fuente: elaboración propia.}
Terreno & 6 apoyos × 7 turnos × 12 h & Apoyos temporales de OPS/TEC & 504 \\
Control de marcha & 40 HH × 3 semanas & LIG 36 + ADT adicional 84 & 120 \\
Coordinación & 8 HH × 3 semanas & LIG 12 + LF 12 & 24 \\
Ajustes de interfaz & DEV 56 + QA 8 & Apoyos adicionales & 64 \\
Capacitación nocturna & 16 sesiones × 2 HH de relator & Relator adicional; hasta 16 asistentes por sesión & 32 \\
Acreditación & 2 apoyos fuera de la cobertura de R-23 × 8 HH & Apoyo de coordinación & 16 \\
Total bruto & 504 + 120 + 24 + 64 + 32 + 16 & 748 HH incrementales + 12 HH reordenadas & 760 \\
\end{tabla}
LIG/LDA suma 94,33 HH base en M16; con 48 HH de este evento llega a 142,33 de 152 HH nominales, margen 9,67. LF reordena 12 HH de preparación anticipada de la primera priorización evolutiva hacia M21 dentro de 1.2.5.4: se reservan allí 12 de sus 152 HH continuas, sin añadir 12 sobre esa capacidad ni recortar las revisiones y aceptaciones E2. Los 64 HH de ajustes usan apoyos adicionales y no vuelven a consumir las 1.280 HH de corrección. El plan del evento no acredita una contratación ya ejecutada.

Las 32 HH de relator cubren hasta 256 asistentes pendientes con sesiones de dos horas y grupos de hasta dieciséis: 2 × techo(asistentes pendientes / 16). Si la cohorte supera ese límite, se amplía la relatoría. Las dos acreditaciones adicionales son distintas de las cuatro previstas en R-23; si estas cuatro no pueden reutilizarse para las seis posiciones de terreno, se agregan otras 32 HH de acreditación. Los escenarios de varias áreas comparten control y ajustes sólo mientras no cambie esa hipótesis de alcance.

La extensión residual de dos semanas estima 528 HH brutas y 520 incrementales; conserva 1.130 HH como valoración residual y 226 HH de valor esperado. Dos áreas durante tres semanas requieren 1.312 HH brutas; tres requieren 1.864 y cuatro 2.416. Desde tres áreas, o ante otro cambio de turnos/duración, se aprueba capacidad y financiación adicional antes de activar el plan. Si se materializa el escenario de un área, las 748 HH incrementales superan sus 226 HH esperadas en 522 HH: esa diferencia es una solicitud de consumo adicional de contingencia o gestión, no una reasignación de gestión ya aprobada. Las fechas contractuales y las condiciones de Art. 17.3 se conservan.

\seccion{Recuperación y alternativa de continuidad}
Se adopta cuatro horas para restablecer el registro operacional ante pérdida total de sala, conforme a S4 4.2.3.4, S5 Tabla \refext{S5}{tab:objetivos-recuperacion-mecanismo-datos-sostiene} y Art. 20. Se prepara el kit fuera del recinto, se verifica la sala de 2014 y se conservan hechos posteriores a la copia mediante retención o réplica exterior y reproducción sin duplicados. La Tabla \ref{tab:presupuesto-recuperacion-recinto} es un presupuesto de diseño, no el resultado de una prueba. La copia horaria aislada no acredita por sí sola cero eventos perdidos ni RPO de quince minutos.

\begin{tabla}{I{0.360}I{0.120}I{0.520}}{Presupuesto de recuperación del recinto.}{tab:presupuesto-recuperacion-recinto}
  {\cab{Paso} & \cab{Minutos} & \cab{Evidencia}}
  {Fuente: elaboración propia.}
Detectar y activar & 15 & Inicio y responsable trazables \\
Kit y rutas preconfiguradas & 45 & Cómputo y comunicación disponibles \\
Desplegar configuración y restaurar base & 60 & Nodo iniciado con identidades válidas \\
Reproducir hechos y reconciliar & 60 & Eventos completos, sin duplicados \\
Probar transacciones críticas & 30 & Restablecimiento de extremo a extremo \\
Margen de ejecución & 30 & Variación absorbida \\
Total & 240 & Objetivo máximo de 4 horas \\
\end{tabla}
Se programan 96 HH nuevas: preparación de 32 HH en M6–M8 (AS 16 + SRE 8 + TEC 8), y dos ejercicios de 32 HH en M15/M21 (cada uno SRE 8 + OPS 8 + TEC 16). Se mantienen separadas las 160 HH vigentes de corte de enlace y falla de servidor. Los ejercicios simulan indisponibilidad íntegra de sala, incluyen desconexión de nube y concilian eventos; sus informes y brechas cerradas son evidencias de H7/H12. Las 96 HH excluyen adquisiciones, traslados e instalación del kit y adecuaciones de sala, que se dimensionan en el plan de adquisiciones.

Para R-20, Brazil South sigue siendo el sitio del diseño vigente y requiere autorización previa escrita antes de la primera transferencia. Las 40 HH preventivas preparan el expediente alternativo (AS 24 + BCP externo 16). Si se activa una adaptación chilena de plataforma compatible, la Tabla \ref{tab:adaptacion-condicionada-recuperacion} desglosa sus 400 HH incrementales. La aprobación y la prueba real anteceden a producción; la ubicación en Chile no resuelve por sí sola los tratamientos de servicios globales.

\begin{tabla}{I{0.240}I{0.360}I{0.080}I{0.320}}{Adaptación condicionada de la recuperación.}{tab:adaptacion-condicionada-recuperacion}
  {\cab{Trabajo} & \cab{Cálculo por rol} & \cab{HH} & \cab{Paquete ejecutor}}
  {Fuente: elaboración propia.}
Cumplimiento y contratación & ESI 16 + ABG 16 & 32 & 1.3.6.1–1.3.6.2; 1.1.6.2 \\
Diseño y amenazas & AS 40 + BCP 24 & 64 & 1.3.1.4 \\
Infraestructura y controles & AS 24 + DEV 64 + ESI 24 & 112 & 1.6.1.1 \\
Réplica y reconciliación & DEV 64 + SRE 24 + ESI 8 & 96 & 1.6.1.1 \\
Conmutación y retorno & SRE 24 + OPS 24 + TEC 16 & 64 & 1.5.5.1 \\
Manuales y transferencia & SRE 16 + OPS 8 + AS 8 & 32 & 1.6.5.2–1.6.5.3 \\
Total & 32 + 64 + 112 + 96 + 64 + 32 & 400 & Activación condicionada de R-20 \\
\end{tabla}
Se propone diseño en M5 y habilitación M5–M8 según contratación, con ensayo antes del hito afectado. Los incrementos se cubren con dedicación parcial disponible o apoyos adicionales y se nivelan antes de ejecutar. Si exige sustituir plataformas, construir un centro o cambiar servicios globales, se estima otro alcance. Las 400 HH cubren adaptación, no consumo de infraestructura ni adquisiciones: se presupuesta el incremento neto contra el sitio sustituido en la Oferta Económica. El escenario activado suma 440 HH adicionales, 40 preventivas y 400 contingentes; no se suma el VE de 40 HH a las 400 HH del mismo evento.

\seccion{Evidencias y referencias de las decisiones}
El Comunicado 10 exige cifras reproducibles, coherencia de diseño y trazabilidad entre requisitos, implementación y operación (CMAA, 2026d, secciones 3 y 7.1). Por ello se conserva el cálculo exacto de R-15, se muestra la dimensión dominante de impacto y se separan objetivos, estimaciones y resultados de ensayos. Las decisiones no modifican los objetivos contractuales ni declaran autorizaciones recibidas. Los informes de prueba, actas de cierre y contratos de recursos se registran en T-13/T-17 y en los paquetes ejecutores.

\seccionsinnumero{Referencias}
\begin{referencias}
\item Compañía Minera Altos de Aranda. (2026a). \emph{Bases administrativas}. Licitación N° TFEP-01/2026.
\item Compañía Minera Altos de Aranda. (2026b). \emph{Bases técnicas del caso Minería}. Licitación N° TFEP-01/2026.
\item Compañía Minera Altos de Aranda. (2026c). \emph{Bases técnicas transversales}. Licitación N° TFEP-01/2026.
\item Compañía Minera Altos de Aranda. (2026d, 23 de septiembre). \emph{Comunicado 10: Estructura obligatoria de las propuestas preparatorias y técnica final} [Comunicación de licitación]. Licitación N° TFEP-01/2026.
\item Overfore S.A. (2026a). Formulario T-14: Plan de trabajo, EDT y carta Gantt. Licitación N° TFEP-01/2026.
\item Overfore S.A. (2026b). Formulario T-16: Plan de riesgos. Licitación N° TFEP-01/2026.
\item Overfore S.A. (2026c). Anexos del Subdocumento 8: Plan de riesgos. Licitación N° TFEP-01/2026.
\item Overfore S.A. (2026d). Subdocumento 4 Arquitectura lógica y física de la solución. Licitación N° TFEP-01/2026.
\item Overfore S.A. (2026e). Subdocumento 5 Modelo y gestión de datos. Licitación N° TFEP-01/2026.
\item Overfore S.A. (2026f). Subdocumento 7 Plan de trabajo, EDT, cronograma e implantación. Licitación N° TFEP-01/2026.
\item Overfore S.A. (2026g). Formulario T-15: Nivelación de recursos. Licitación N° TFEP-01/2026.
\end{referencias}
\end{horizontal}
\end{document}
